\label{sec:background}

\subsection{NLCD 2021: Data and Coverage}

The National Land Cover Database (NLCD) is produced by the Multi-Resolution
Land Characteristics (MRLC) Consortium, a partnership of federal agencies
including the USGS, EPA, USDA Forest Service, and National Park Service
\citep{homer2020nlcd}.
NLCD provides a consistent, wall-to-wall classification of land cover for
the contiguous United States at 30-meter spatial resolution, derived from
Landsat 8 satellite imagery supplemented by ancillary datasets (building
footprints, impervious surface maps, canopy cover estimates).

The 2021 vintage is the most recent NLCD release as of 2024.
It is available as a single GeoTIFF covering CONUS (approximately 2~GB
compressed), in Albers Equal Area projection (EPSG:5070), with 30-meter
pixel resolution.
The data is published at \url{https://www.mrlc.gov/data/nlcd-2021-land-cover-conus}
and is public domain under U.S.\ government work provisions.

\paragraph{NLCD class hierarchy.}
NLCD uses a hierarchical 16-category system derived from the Anderson
Level II classification \citep{anderson1976land}.
Table~\ref{tab:nlcd-classes} lists all classes with their integer codes
and the M.2 category assignment.

\begin{table}[ht]
\centering
\caption{NLCD 2021 land cover classes and M.2 category assignments.
  Class~11 is excluded from the cosine similarity vector and used only
  for the hard-cut rule.}
\label{tab:nlcd-classes}
\small
\begin{tabular}{rll}
\toprule
Code & NLCD class name & M.2 category \\
\midrule
11   & Open Water & Hard-cut trigger (not in vector) \\
21   & Developed, Open Space & \texttt{pct\_open} \\
22   & Developed, Low Intensity & \texttt{pct\_residential} \\
23   & Developed, Medium Intensity & \texttt{pct\_residential} \\
24   & Developed, High Intensity & \texttt{pct\_commercial} \\
31   & Barren Land & \texttt{pct\_natural} \\
41   & Deciduous Forest & \texttt{pct\_natural} \\
42   & Evergreen Forest & \texttt{pct\_natural} \\
43   & Mixed Forest & \texttt{pct\_natural} \\
52   & Shrub/Scrub & \texttt{pct\_natural} \\
71   & Grassland/Herbaceous & \texttt{pct\_natural} \\
81   & Pasture/Hay & \texttt{pct\_natural} \\
82   & Cultivated Crops & \texttt{pct\_natural} \\
90   & Woody Wetlands & \texttt{pct\_natural} \\
95   & Emergent Herbaceous Wetlands & \texttt{pct\_natural} \\
\bottomrule
\end{tabular}
\end{table}

The NLCD ``developed'' classes (21--24) form a gradient of impervious surface
density.
Class~21 (developed, open space, $<$20\% impervious) includes suburban lawns,
golf courses, parks in developed areas, and cemeteries.
Class~22 (low intensity, 20--49\% impervious) is the predominant classification
for single-family residential neighborhoods.
Class~23 (medium intensity, 50--79\% impervious) captures denser residential
development: row houses, townhouses, multi-family housing, and mixed-use
residential buildings.
Class~24 (high intensity, $\geq$80\% impervious) captures commercial districts,
industrial parks, transportation hubs, and urban core blocks where nearly
the entire land surface is covered by buildings, pavement, or other
impervious materials.

\paragraph{Zonal statistics computation.}
NLCD is a raster product; census tracts are polygon features from TIGER/Line
shapefiles \citep{census2020tiger}.
Computing land use character vectors requires zonal statistics: for each tract
polygon, count the number of 30-meter pixels of each NLCD class within the
polygon boundary.

The computation proceeds as follows:
\begin{enumerate}
  \item Reproject TIGER/Line tract polygons from EPSG:4326 to EPSG:5070
    (NLCD native projection) to avoid reprojection artifacts in the raster
    sampling.
  \item For each tract, use the GDAL \texttt{zonal\_stats} function to count
    pixels within the polygon boundary (all-touched: false; pixel center
    must fall within polygon).
  \item Aggregate pixel counts by M.2 category to produce five per-tract
    counts: \texttt{n\_residential}, \texttt{n\_commercial},
    \texttt{n\_open}, \texttt{n\_natural}, \texttt{n\_water}.
  \item Divide by total pixel count to produce fractions.
\end{enumerate}

\paragraph{Edge cases.}
Census tracts in coastal areas or along large rivers may have substantial water
pixels (class~11) within their polygons even if the residential area of the
tract is inland.
The water fraction is excluded from the cosine similarity vector to prevent
high water pixel counts from suppressing similarity between two inland
residential tracts that happen to border the same body of water.
Instead, water is handled by the hard-cut rule described in
Section~\ref{sec:algorithm}.

\subsection{Physical Geography as Communities of Interest}

The use of physical geography as a redistricting criterion predates the modern
communities-of-interest doctrine.
The first American redistricting practitioners drew district lines along rivers,
mountain ridges, and coastal features because those were the natural divisions
visible on maps and meaningful to the communities they were drawing.

\citet{schwartzberg1966reapportionment} noted the relationship between
geographic compactness and natural land use boundaries: compact districts tend
to follow natural divisions, and natural divisions tend to reflect the physical
landscape that shapes community identity.
The connection is not merely aesthetic.
Communities that share a physical landscape share infrastructure needs
(flood management, shoreline erosion, wildfire risk, agricultural water
management) that translate directly into shared legislative interests.

\paragraph{Urban-rural transitions.}
The transition from residential suburban development to commercial corridor,
or from developed land to agricultural land, is one of the most consistent
physical geographic community boundaries in American metropolitan areas.
Residents of a suburban neighborhood and workers in an adjacent industrial
park may share a common boundary on a map but live in economically and
physically distinct communities.
The NLCD classification makes this distinction explicit: one side of the
boundary is class~22 (low-intensity residential), the other is class~24
(high-intensity developed, commercial/industrial).

\paragraph{Water bodies as community barriers.}
Rivers, bays, sounds, and lakes create natural community boundaries by
imposing travel barriers and defining watershed communities with shared
environmental concerns.
North Carolina's Pamlico Sound divides the coastal mainland from the
Outer Banks, creating communities whose access, infrastructure, and
economic character are fundamentally different.
Louisiana's bayou network defines communities more precisely than parish
boundaries.
Texas's Gulf Coast bays and barrier islands define distinct communities
not well captured by county lines.
The hard-cut rule in M.2 encodes this physical reality as a zero-weight
edge: communities separated by majority-water adjacencies are prevented
from being placed in the same district unless explicit bridge adjacency
is recorded in the graph structure.
